\section{Scalar dense inner refinements}
\label{sec:scalar-dense-inner-refinements}

This note preserves scalar dense-inner refinements that are useful for finite diagnostics and possible sharper theorem
work, but are not part of the current compiled main proof spine. The main paper now keeps only the lemmas needed for the
random sliding dense asymptotic theorem: the three-term tail, run lower-tail exponent, forced-termination contraction,
and the packaged inner interface.

\subsection{Run-tail exponent derivation}

This note preserves the expanded Stirling calculation behind the run lower-tail exponent used in the main scalar dense
inner section.  Let \(U\) be uniform over weight-\(w\) words in \(\{0,1\}^N\), set
\[
w=\lfloor\omega N\rfloor,\qquad r=\lfloor\rho N\rfloor,
\]
and assume \(0<\rho<\omega(1-\omega)\).  The exact run-count formula gives
\[
\Pr[R(U)<r]
=
\frac{1}{\binom{N}{w}}
\sum_{s=1}^{r-1}
\binom{w-1}{s-1}\binom{N-w+1}{s}.
\]
For \(s=\lfloor\sigma N\rfloor\), Stirling's formula gives
\[
\log_2\binom{w-1}{s-1}
=
\omega N H_2\!\left(\frac{\sigma}{\omega}\right)+O(\log N),
\]
\[
\log_2\binom{N-w+1}{s}
=
(1-\omega)N H_2\!\left(\frac{\sigma}{1-\omega}\right)+O(\log N),
\]
and
\[
\log_2\binom{N}{w}=NH_2(\omega)+O(\log N).
\]
Thus the summand exponent is
\[
-N\Psi_{\rm run}(\omega,\sigma)+O(\log N),
\]
where
\[
\Psi_{\rm run}(\omega,\sigma)
:=
H_2(\omega)
-\omega H_2\!\left(\frac{\sigma}{\omega}\right)
-(1-\omega)H_2\!\left(\frac{\sigma}{1-\omega}\right).
\]
The derivative is
\[
\frac{\partial}{\partial\sigma}\Psi_{\rm run}(\omega,\sigma)
=
\log_2\!\left(\frac{\sigma^2}{(\omega-\sigma)(1-\omega-\sigma)}\right),
\]
which is negative for \(\sigma<\omega(1-\omega)\).  Hence the largest summand in the lower tail occurs at the boundary
\(\sigma=\rho+O(1/N)\), and the whole sum is bounded by
\[
N^{O(1)}2^{-N\Psi_{\rm run}(\omega,\rho)}.
\]

\subsection{Gap, isolated-slice, and paired-cover refinements}
\begin{lemma}[Gap composition conditioned on weight and run count]
\label{lem:gap-composition-given-runs}
Let $U$ be uniform over $\{0,1\}^N$ conditioned on $\wt(U)=w\ge 1$ and $R(U)=s$.
Write the zero gaps around the $s$ runs as
\[
z_0,\ z_1,\ \dots,\ z_{s-1},\ z_s,
\]
where $z_0$ and $z_s$ are the leading and trailing zero counts and each internal gap $z_i$ for
$1\le i\le s-1$ is at least $1$. Define shifted internal gaps
\[
y_0:=z_0,\qquad y_i:=z_i-1 \ (1\le i\le s-1),\qquad y_s:=z_s.
\]
Then $(y_0,\dots,y_s)$ is uniform over all weak compositions of
\[
M:=N-w-(s-1)
\]
into $s+1$ parts.

In particular, for any fixed internal gap $G=z_i$ with $1\le i\le s-1$,
\begin{equation}
\label{eq:internal-gap-marginal}
\PP[G=g \mid \wt(U)=w,\ R(U)=s]
\;=\;
\frac{\binom{M-g+s-1}{s-1}}{\binom{M+s}{s}}
\qquad (g\ge 1),
\end{equation}
whenever $M-g\ge 0$, and therefore
\begin{equation}
\label{eq:internal-gap-mean}
\EE[G \mid \wt(U)=w,\ R(U)=s]
\;=\;
1+\frac{M}{s+1}
\;=\;
\frac{N-w+2}{s+1}.
\end{equation}
\end{lemma}

\begin{proof}
Conditioned on $\wt(U)=w$ and $R(U)=s$, a word is specified by:
\begin{itemize}
\item a composition of $w$ into $s$ positive run lengths, and
\item a choice of zero gaps $z_0,z_1,\dots,z_s$ with $z_0,z_s\ge 0$, $z_i\ge 1$ for $1\le i\le s-1$, and
\(\sum_{i=0}^s z_i=N-w\).
\end{itemize}
After the shift $y_i=z_i-1$ on the internal gaps, the second datum becomes a weak composition of
$M=N-w-(s-1)$ into $s+1$ parts. For each fixed shifted gap vector there are exactly
$\binom{w-1}{s-1}$ compatible run-length compositions, so under the uniform law on the $(w,s)$-slice the shifted
gap vector is itself uniform over those weak compositions.

Fix an internal gap $G=z_i$ and write $G=1+Y_i$. The event $\{G=g\}$ is equivalent to $Y_i=g-1$, leaving
$M-(g-1)=M-g+1$ units to be distributed across the remaining $s$ coordinates. Hence
\[
\PP[G=g \mid \wt(U)=w,\ R(U)=s]
\;=\;
\frac{\binom{M-g+s-1}{s-1}}{\binom{M+s}{s}},
\]
which is \eqref{eq:internal-gap-marginal}. Finally, by symmetry of the uniform composition law,
$\EE[Y_i]=M/(s+1)$, so \eqref{eq:internal-gap-mean} follows from $G=1+Y_i$.
\end{proof}

\begin{corollary}[Forced-cluster count from short gaps]
\label{cor:forced-cluster-count}
Fix a memory parameter $m\ge 1$. Under the same conditioning as
Lemma~\ref{lem:gap-composition-given-runs}, define the \emph{forced cluster count}
\[
C_m(U)
\;:=\;
1+\#\{\,i\in\{1,\dots,s-1\}: z_i\ge m\,\}.
\]
Equivalently, one merges adjacent runs whenever their separating zero gap is $<m$.
Then
\begin{equation}
\label{eq:forced-cluster-prob}
\PP[z_i\ge m \mid \wt(U)=w,\ R(U)=s]
\;=\;
\frac{\binom{M-m+s+1}{s}}{\binom{M+s}{s}}
\end{equation}
for each internal gap \(z_i\), whenever \(M-m+1\ge 0\), and therefore
\begin{equation}
\label{eq:forced-cluster-mean}
\EE[C_m(U)\mid \wt(U)=w,\ R(U)=s]
\;=\;
1+(s-1)\frac{\binom{M-m+s+1}{s}}{\binom{M+s}{s}}.
\end{equation}
\end{corollary}

\begin{proof}
By definition, \(z_i\ge m\) if and only if \(y_i=z_i-1\ge m-1\).
Under Lemma~\ref{lem:gap-composition-given-runs}, the shifted gap vector is uniform over weak compositions of
\(M\) into \(s+1\) parts. So fixing \(y_i\ge m-1\) and subtracting \(m-1\) from that coordinate leaves a weak
composition of \(M-(m-1)\) into \(s+1\) parts. Hence
\[
\PP[z_i\ge m \mid \wt(U)=w,\ R(U)=s]
\;=\;
\frac{\binom{M-(m-1)+s}{s}}{\binom{M+s}{s}}
\;=\;
\frac{\binom{M-m+s+1}{s}}{\binom{M+s}{s}},
\]
which is \eqref{eq:forced-cluster-prob}. The expectation formula
\eqref{eq:forced-cluster-mean} then follows by linearity of expectation over the \(s-1\) exchangeable internal gaps.
\end{proof}

\begin{lemma}[First-run late tail conditioned on the run count]
\label{lem:first-run-late-given-runs}
Let $U$ be uniform over $\{0,1\}^N$ conditioned on $\wt(U)=w\ge 1$ and $R(U)=s$.
Let $T_1(U)$ denote the start position of the first run of ones.
Then for every integer $t\in\{0,1,\dots,N-w\}$,
\begin{equation}
\label{eq:first-run-late-given-runs}
\PP[T_1(U)>t \mid \wt(U)=w,\ R(U)=s]
\;=\;
\frac{\binom{N-t-w+1}{s}}{\binom{N-w+1}{s}}.
\end{equation}
Consequently,
\begin{equation}
\label{eq:first-run-late-given-runs-simple}
\PP[T_1(U)>t \mid \wt(U)=w,\ R(U)=s]
\;\le\;
\left(1-\frac{t}{N}\right)^s.
\end{equation}
\end{lemma}

\begin{proof}
The event $\{T_1(U)>t\}$ means that the first $t$ coordinates are all zero, so the entire weight-$w$, $s$-run pattern
must be placed in the remaining $N-t$ coordinates. By Lemma~\ref{lem:run-count-distribution}, the number of such words is
\[
\binom{w-1}{s-1}\binom{N-t-w+1}{s},
\]
while the total number of weight-$w$, $s$-run words is
\[
\binom{w-1}{s-1}\binom{N-w+1}{s}.
\]
Taking the ratio gives \eqref{eq:first-run-late-given-runs}.

For the bound, write
\[
\frac{\binom{N-t-w+1}{s}}{\binom{N-w+1}{s}}
\;=\;
\prod_{j=0}^{s-1}\frac{N-t-w+1-j}{N-w+1-j}
\;=\;
\prod_{j=0}^{s-1}\left(1-\frac{t}{N-w+1-j}\right).
\]
Since $N-w+1-j\le N$, each factor is at most $1-t/N$, proving
\eqref{eq:first-run-late-given-runs-simple}.
\end{proof}

\begin{corollary}[Conditioned late-start base factor]
\label{cor:first-run-late-qbase}
Fix $\delta\in(0,1/2)$ and let
\[
L:=\lceil 2\delta N\rceil.
\]
Under the same conditioning as Lemma~\ref{lem:first-run-late-given-runs},
\begin{equation}
\label{eq:first-run-late-qbase}
\PP[T_1(U)>N-L \mid \wt(U)=w,\ R(U)=s]
\;=\;
\frac{\binom{L-w+1}{s}}{\binom{N-w+1}{s}}
\;\le\;
\left(\frac{L}{N}\right)^s.
\end{equation}
In particular,
\begin{equation}
\label{eq:first-run-late-qbase-asym}
\PP[T_1(U)>N-L \mid \wt(U)=w,\ R(U)=s]
\;\le\;
\left(2\delta + O\!\left(\frac{1}{N}\right)\right)^s.
\end{equation}
\end{corollary}

\begin{proof}
Apply Lemma~\ref{lem:first-run-late-given-runs} with $t=N-L$.
Then
\[
\PP[T_1(U)>N-L \mid \wt(U)=w,\ R(U)=s]
\;=\;
\frac{\binom{L-w+1}{s}}{\binom{N-w+1}{s}}.
\]
The bound follows from
\[
\frac{\binom{L-w+1}{s}}{\binom{N-w+1}{s}}
\;=\;
\prod_{j=0}^{s-1}\frac{L-w+1-j}{N-w+1-j}
\;\le\;
\left(\frac{L}{N}\right)^s,
\]
and $L/N = 2\delta + O(1/N)$ gives \eqref{eq:first-run-late-qbase-asym}.
\end{proof}

\begin{corollary}[Two isolated ones: exact late-placement factor]
\label{cor:first-run-late-two-run}
Fix $\delta\in(0,1/2)$ and let $L=\lceil 2\delta N\rceil$.
If $U$ is uniform over $\{0,1\}^N$ conditioned on $\wt(U)=2$ and $R(U)=2$, then
\begin{equation}
\label{eq:first-run-late-two-run}
\PP[T_1(U)>N-L \mid \wt(U)=2,\ R(U)=2]
\;=\;
\frac{\binom{L-1}{2}}{\binom{N-1}{2}}.
\end{equation}
Equivalently,
\begin{equation}
\label{eq:first-run-late-two-run-asym}
\PP[T_1(U)>N-L \mid \wt(U)=2,\ R(U)=2]
\;=\;
\left(\frac{L}{N}\right)^2
\left(1+O\!\left(\frac{1}{L}\right)\right).
\end{equation}
\end{corollary}

\begin{proof}
This is the special case $(w,s)=(2,2)$ of Corollary~\ref{cor:first-run-late-qbase}. The asymptotic form follows from
\[
\frac{\binom{L-1}{2}}{\binom{N-1}{2}}
\;=\;
\left(\frac{L}{N}\right)^2
\frac{(1-1/L)(1-2/L)}{(1-1/N)(1-2/N)}.
\]
\end{proof}

\begin{corollary}[Isolated order statistics: exact late tail]
\label{cor:isolated-order-late}
Fix integers \(1\le r\le h\le N\) and let \(L\in\{1,\dots,N\}\).
Let \(U\) be uniform over \(\{0,1\}^N\) conditioned on \(\wt(U)=h\) and \(R(U)=h\), and write
\[
X_1(U)<X_2(U)<\cdots<X_h(U)
\]
for the ordered one positions.
Then
\begin{equation}
\label{eq:isolated-order-late}
\PP[X_r(U)>N-L \mid \wt(U)=h,\ R(U)=h]
\;=\;
\frac{1}{\binom{N-h+1}{h}}
\sum_{j=0}^{r-1}
\binom{N-L-r+1}{j}\binom{L-h+r}{h-j},
\end{equation}
with the convention that the sum is interpreted as \(1\) when \(N-L-r+1<0\).
\end{corollary}

\begin{proof}
Conditioned on \(\wt(U)=h\) and \(R(U)=h\), the isolated one positions satisfy
\[
1\le X_1<\cdots<X_h\le N,
\qquad
X_{i+1}\ge X_i+2.
\]
Define
\[
Y_i:=X_i-i+1,\qquad i=1,\dots,h.
\]
Then
\[
1\le Y_1<\cdots<Y_h\le N-h+1,
\]
and the map \(X\mapsto Y\) is a bijection from isolated \(h\)-subsets of \([N]\) to ordinary \(h\)-subsets of
\([N-h+1]\). Hence \(Y_1<\cdots<Y_h\) is uniform over all \(h\)-subsets of \([N-h+1]\).

Now \(X_r>N-L\) is equivalent to
\[
Y_r=X_r-r+1 > N-L-r+1.
\]
So among the \(Y_i\), at most \(r-1\) may lie in the prefix \(\{1,\dots,N-L-r+1\}\), while the remaining
\(h-j\) elements must lie in the tail of size
\[
(N-h+1)-(N-L-r+1)=L-h+r.
\]
Choosing exactly \(j\) prefix elements gives
\[
\binom{N-L-r+1}{j}\binom{L-h+r}{h-j},
\]
and summing over \(j=0,\dots,r-1\) proves \eqref{eq:isolated-order-late}.
\end{proof}

\begin{corollary}[Isolated early-count decomposition]
\label{cor:isolated-early-count}
Fix \(h\le N\) and let \(L\in\{1,\dots,N\}\).
Let \(U\) be uniform over \(\{0,1\}^N\) conditioned on \(\wt(U)=h\) and \(R(U)=h\), and define
\[
J_L(U)
\;:=\;
\#\{\,i\in\{1,\dots,h\}: X_i(U)\le N-L\,\},
\]
the number of isolated one-positions that start before the final \(L\) coordinates.
Then for every \(j\in\{0,\dots,h\}\),
\begin{equation}
\label{eq:isolated-early-count}
\PP[J_L(U)=j \mid \wt(U)=h,\ R(U)=h]
\;=\;
\frac{\binom{N-L}{j}\binom{L-h+1}{h-j}}{\binom{N-h+1}{h}},
\end{equation}
with the convention that the right-hand side is \(0\) when either binomial coefficient is undefined.
\end{corollary}

\begin{proof}
Under the shift \(Y_i=X_i-i+1\) from the proof of Corollary~\ref{cor:isolated-order-late},
the isolated slice \(\wt(U)=h,\ R(U)=h\) becomes a uniform \(h\)-subset
\[
1\le Y_1<\cdots<Y_h\le N-h+1.
\]
The condition \(X_i\le N-L\) is equivalent to \(Y_i\le N-L-i+1\), and therefore
\(J_L(U)=j\) means exactly \(j\) of the shifted positions lie in the prefix \(\{1,\dots,N-L\}\),
while the remaining \(h-j\) lie in the tail of size
\[
(N-h+1)-(N-L)=L-h+1.
\]
Choosing those two groups independently gives
\[
\binom{N-L}{j}\binom{L-h+1}{h-j}
\]
valid isolated \(h\)-subsets. Dividing by the total \(\binom{N-h+1}{h}\) proves
\eqref{eq:isolated-early-count}.
\end{proof}

\begin{corollary}[Expected number of early isolated pairs]
\label{cor:isolated-early-pair-mean}
Under the same conditioning and notation as Corollary~\ref{cor:isolated-early-count},
\begin{equation}
\label{eq:isolated-early-pair-mean}
\EE\!\left[\binom{J_L(U)}{2}\,\middle|\, \wt(U)=h,\ R(U)=h\right]
\;=\;
\binom{h}{2}\,
\frac{\binom{N-L}{2}}{\binom{N-h+1}{2}}.
\end{equation}
\end{corollary}

\begin{proof}
Under the shift \(Y_i=X_i-i+1\) from the proof of Corollary~\ref{cor:isolated-order-late}, the isolated slice becomes a
uniform \(h\)-subset of \([N-h+1]\), and \(J_L(U)\) counts how many of those \(h\) points land in the distinguished
prefix of size \(N-L\). So \(J_L(U)\) has the usual hypergeometric distribution with population size \(N-h+1\),
success population \(N-L\), and sample size \(h\). The factorial second moment of a hypergeometric variable gives
\[
\EE[J_L(U)(J_L(U)-1)\mid \wt(U)=h,\ R(U)=h]
\;=\;
h(h-1)\frac{(N-L)(N-L-1)}{(N-h+1)(N-h)}.
\]
Dividing by \(2\) proves \eqref{eq:isolated-early-pair-mean}.
\end{proof}

\begin{corollary}[Tiny weights are almost all isolated]
\label{cor:run-count-full}
Let $U$ be uniform over $\{0,1\}^N$ conditioned on $\wt(U)=h\ge 1$.
Then
\begin{equation}
\label{eq:run-count-full}
\PP[R(U)=h]
\;=\;
\frac{\binom{N-h+1}{h}}{\binom{N}{h}}.
\end{equation}
In particular, for every fixed $h$,
\begin{equation}
\label{eq:run-count-full-asym}
\PP[R(U)=h]
\;=\;
1-O_h\!\left(\frac{1}{N}\right),
\end{equation}
and more concretely,
\begin{equation}
\label{eq:run-count-full-lower}
\PP[R(U)=h]
\;\ge\;
1-\frac{h(h-1)}{N}.
\end{equation}
\end{corollary}

\begin{proof}
The exact formula is the case $s=h$ of Lemma~\ref{lem:run-count-distribution}. For the lower bound, let
\[
A_i:=\{U_i=U_{i+1}=1\},\qquad i=1,\dots,N-1.
\]
If $R(U)<h$, then at least one adjacency event $A_i$ occurs. Under the uniform weight-$h$ model,
\[
\PP[A_i]
\;=\;
\frac{\binom{N-2}{h-2}}{\binom{N}{h}}
\;=\;
\frac{h(h-1)}{N(N-1)}.
\]
Hence by a union bound,
\[
\PP[R(U)<h]
\;\le\;
(N-1)\frac{h(h-1)}{N(N-1)}
\;=\;
\frac{h(h-1)}{N},
\]
which proves \eqref{eq:run-count-full-lower}. The asymptotic claim \eqref{eq:run-count-full-asym} is immediate for
fixed \(h\).
\end{proof}

\begin{lemma}[Fixed-tap late-start plus one-episode correction]
\label{lem:dense-late-start-one-episode}
Let $U$ be uniform over $\{0,1\}^N$ conditioned on $\wt(U)=w\ge 1$ and $R(U)=s$.
Let $Y=\Enc_{\mathrm{in}}(U)$, fix an integer $L$ with $1\le L\le N$, and let $T_1(U)$ be the first run start.
For a fixed input \(u\), let
\[
B_{T_1}(u;L):=
|\{j\in\supp(u):T_1(u)<j\le T_1(u)+L\}|.
\]
Then for every threshold $d\ge 0$,
\begin{equation}
\label{eq:dense-late-start-one-episode}
\PP[\wt(Y)\le d \mid \wt(U)=w,\ R(U)=s]
\;\le\;
\PP[T_1(U)>N-L \mid \wt(U)=w,\ R(U)=s]
\;+\;
2^{-(m-1)}\EE[B_{T_1}(U;L)\mid \wt(U)=w,\ R(U)=s]
\;+\;
\PP[\Binomial(L,\tfrac12)\le d].
\end{equation}
\end{lemma}

\begin{proof}
Let $E_{\mathrm{late}}:=\{T_1(U)>N-L\}$.
On the complementary event $E_{\mathrm{late}}^c$, the first run starts by time $N-L$.
As in the proof of Corollary~\ref{cor:dense-inner-single-run}, that first run forces the encoder into the ON state at
the next time step, thereby initiating the first ON-episode.

If that first ON-episode terminates within its first $L$ steps, then
Lemma~\ref{lem:fixedtap-one-episode-termination} bounds its probability by
\(B_{T_1}(U;L)2^{-(m-1)}\) for the fixed input \(U\).

Otherwise the first ON-episode survives for at least $L$ steps. Since the total weight of $Y$ is at most $d$, in
particular those first \(L\) output bits of the episode must contain at most \(d\) ones.
Lemma~\ref{lem:fixedtap-surviving-episode-tail} bounds the joint probability of that event by
\(\PP[\Binomial(L,\tfrac12)\le d]\).

Thus, on $E_{\mathrm{late}}^c$, the low-weight event is contained in the union of the ``early termination within $L$''
event and the length-$L$ binomial tail event. Adding back the late-start event $E_{\mathrm{late}}$ proves
\eqref{eq:dense-late-start-one-episode}.
\end{proof}

\begin{corollary}[Run-mixture slice bound from one episode]
\label{cor:dense-late-start-slice}
Fix $\delta\in(0,1/2)$ and let $d=\lfloor \delta N\rfloor$.
Then for every $w\ge 1$,
\begin{equation}
\label{eq:dense-late-start-slice}
p_w(\delta)
\;\le\;
\sum_{s=1}^{w}
\PP[R(U)=s]\,
\frac{\binom{L-w+1}{s}}{\binom{N-w+1}{s}}
\;+\;
w2^{-(m-1)}
\;+\;
\PP[\Binomial(L,\tfrac12)\le d],
\end{equation}
where \(U\) is uniform on the weight-\(w\) slice and \(L\in\{1,\dots,N\}\) is arbitrary.
\end{corollary}

\begin{proof}
Average Lemma~\ref{lem:dense-late-start-one-episode} over the conditional run-count law from
Lemma~\ref{lem:run-count-distribution}, then substitute Corollary~\ref{cor:first-run-late-qbase} for the
late-start term. The fixed-tap candidate term is at most \(w2^{-(m-1)}\), since there are at most \(w-1\) later
input-one positions after the first run start.
\end{proof}

\begin{corollary}[Isolated fixed-tap one-episode bound]
\label{cor:fixedtap-isolated-one-episode}
Let \(U\) be uniform conditioned on \(\wt(U)=h\) and \(R(U)=h\). For \(1\le L\le N\) and \(d\ge 0\),
\[
\PP[\wt(Y)\le d\mid \wt(U)=h,\ R(U)=h]
\le
\frac{\binom{L-h+1}{h}}{\binom{N-h+1}{h}}
\;+\;
2^{-(m-1)}\binom{h}{2}\frac{2L}{N-h}
\;+\;
\PP[\Binomial(L,\tfrac12)\le d],
\]
with the convention that undefined binomial coefficients are zero.
\end{corollary}

\begin{proof}
Apply Lemma~\ref{lem:dense-late-start-one-episode} with \(w=s=h\). The late-start term is the exact isolated
late-placement base from Corollary~\ref{cor:first-run-late-qbase}, and the fixed-tap candidate expectation is bounded
by Corollary~\ref{cor:fixedtap-isolated-pair-candidates}.
\end{proof}

\begin{lemma}[Fixed-tap multi-episode early-span certificate]
\label{lem:fixedtap-multiepisode-earlyspan}
Fix an input word \(u\) of weight \(h\), let
\[
E_L(u):=\supp(u)\cap\{1,\dots,N-L\},
\qquad
\Delta_L(u):=
\begin{cases}
0, & E_L(u)=\emptyset,\\
\max E_L(u)-\min E_L(u)+1, & E_L(u)\ne\emptyset,
\end{cases}
\]
and
\[
r_L(u):=\left\lceil \frac{\Delta_L(u)}{L+1}\right\rceil,
\]
and set \(d\ge 0\). Then the fixed-tap dense recursive inner satisfies
\[
\PP[\wt(Y)\le d\mid U=u]
\le
\mathbf{1}_{\{\Delta_L(u)=0\}}
\;+\;
\mathbf{1}_{\{\Delta_L(u)>0\}}\,
\binom{h}{r_L(u)}2^{-r_L(u)(m-1)}
\;+\;
\PP[\Binomial(L,\tfrac12)\le d].
\]
\end{lemma}

\begin{proof}
If \(\Delta_L(u)=0\), the displayed bound is immediate from the first term. Otherwise let \(a:=\min E_L(u)\). At time
\(a\) the encoder is \(\OFF\), so this input one starts an ON episode. If any ON episode that starts at an input-one
position in \(E_L(u)\) survives for \(L\) output times, then Lemma~\ref{lem:fixedtap-surviving-episode-tail} bounds the
joint low-weight event by the displayed binomial tail, choosing the first such surviving episode predictably.

On the complementary event, each ON episode that starts inside the early span terminates within its next \(L\) output
times. A single such episode can cover input-one positions from its start through at most \(L\) later positions, hence
across a span of at most \(L+1\). Since the early input-one positions have span \(\Delta_L(u)\), covering that span
without a surviving length-\(L\) episode requires at least
\(r_L(u)\) genuine fixed-tap terminations. There are at most
\(\binom{h}{r_L(u)}\) possible choices for a subcollection of \(r_L(u)\) terminating input-one positions, and
Lemma~\ref{lem:fixedtap-multicandidate-termination} gives cost \(2^{-r_L(u)(m-1)}\) for each choice. A union bound
proves the claim.
\end{proof}

\begin{corollary}[Full weight-slice multi-episode span bound]
\label{cor:fixedtap-weight-slice-multiepisode}
Let \(h\ge 2\), let \(U\) be uniform on the weight-\(h\) slice, let \(d=\lfloor\delta N\rfloor\), and let
\[
\Pi_{N,h,L}(r):=\PP\!\left[
\left\lceil\frac{\Delta_L(U)}{L+1}\right\rceil=r
\middle| \wt(U)=h
\right].
\]
Then
\[
p_h(\delta)
\le
\Pi_{N,h,L}(0)
\;+\;
\sum_{r\ge 1}
\Pi_{N,h,L}(r)
\left(
\binom{h}{r}2^{-r(m-1)}
\right)
\;+\;
\PP[\Binomial(L,\tfrac12)\le d].
\]
Moreover, the full-slice span law is explicit. Writing \(P=N-L\),
\[
\Pi_{N,h,L}(0)=\frac{\binom{L}{h}}{\binom{N}{h}},
\]
and for \(r\ge 1\),
\[
\Pi_{N,h,L}(r)
=
\frac{1}{\binom{N}{h}}
\sum_{s=(r-1)(L+1)+1}^{\min\{r(L+1),P\}}
(P-s+1)\binom{s-2+L}{h-2},
\]
with the \(s=1\) term interpreted as \(P\binom{L}{h-1}\).
\end{corollary}

\begin{proof}
Average Lemma~\ref{lem:fixedtap-multiepisode-earlyspan} over the uniform weight-\(h\) law. The event
\(\Delta_L(U)=0\) means all \(h\) input ones are in the final \(L\) positions. If the early span is \(s\ge 2\), choose
its first early position in \(P-s+1\) ways, force the endpoints of the span to be selected, and choose the remaining
\(h-2\) input ones from the \(s-2\) interior positions together with the final \(L\) positions. If \(s=1\), choose the
unique early position in \(P\) ways and choose the remaining \(h-1\) positions from the final \(L\) positions. Grouping
the span values by \(r=\lceil s/(L+1)\rceil\) gives the displayed law.
\end{proof}

\begin{definition}[Paired cover of the early support]
\label{def:fixedtap-paired-cover}
Fix an input word \(u\) and \(L\ge 1\). A paired \(L\)-cover of the early support
\[
E_L(u)=\supp(u)\cap\{1,\dots,N-L\}
\]
is a collection of disjoint intervals
\[
I_a=[\alpha_a,\beta_a]\qquad (a=1,\dots,r)
\]
such that
\[
\alpha_a,\beta_a\in\supp(u),\qquad
1\le \beta_a-\alpha_a\le L,
\]
the intervals are ordered by time, and every point of \(E_L(u)\) lies in \(\bigcup_a I_a\).
The total live length of the cover is
\[
T(I_1,\dots,I_r):=\sum_{a=1}^r |I_a|.
\]
\end{definition}

\begin{lemma}[No-survival trajectories induce paired covers]
\label{lem:fixedtap-no-survival-paired-cover}
Fix an input word \(u\) and \(L\ge 1\). On the event that every ON episode whose start position lies in \(E_L(u)\)
terminates within its next \(L\) output times, the early support \(E_L(u)\) admits a paired \(L\)-cover. Moreover the
right endpoint of each cover interval is a genuine fixed-tap termination time.
\end{lemma}

\begin{proof}
Scan the process from left to right. Whenever the encoder is \(\OFF\) at an input-one position
\(\alpha\in E_L(u)\), a new ON episode starts. By the no-survival assumption this episode terminates at some input-one
position \(\beta\le \alpha+L\); Lemma~\ref{lem:fixedtap-termination-needs-input-one} gives that \(\beta\in\supp(u)\).
The interval \([\alpha,\beta]\) covers all input-one positions encountered while that episode is live. After the
termination the encoder is \(\OFF\), so the next uncovered early input-one position, if any, starts the next interval.
The resulting intervals are disjoint, have input-one endpoints, have length at most \(L+1\), and cover \(E_L(u)\).
\end{proof}

\begin{lemma}[Fixed paired-cover stochastic charge]
\label{lem:fixedtap-paired-cover-charge}
Fix an input word \(u\) and a paired \(L\)-cover
\[
I_a=[\alpha_a,\beta_a]\qquad (a=1,\dots,r)
\]
with total live length \(T\). Let \(\mathcal{C}(I_1,\dots,I_r)\) be the event that these intervals are realized as ON
episodes: each \(\alpha_a\) starts an episode, each \(\beta_a\) is the corresponding genuine fixed-tap termination,
and no interval contains an earlier termination. Then for every \(d\ge 0\),
\[
\PP\!\left[
\mathcal{C}(I_1,\dots,I_r)
\wedge
\wt(Y)\le d
\mid U=u
\right]
\le
2^{-r(m-1)}
\Pr\!\left[\Binomial\!\left((T-rm)_+,\tfrac12\right)\le d\right].
\]
\end{lemma}

\begin{proof}
Expose the intervals in time order. For each interval, the fixed-tap termination at its right endpoint costs
\(2^{-(m-1)}\) by the same conditional argument as Lemma~\ref{lem:fixedtap-multicandidate-termination}; multiplying
over the \(r\) endpoints gives the factor \(2^{-r(m-1)}\).

It remains to charge the live outputs not used by the terminal \(m\)-zero windows. In each realized interval the
episode survives up to the terminal zero window. Before that terminal window, every exposed live output either occurs
in the \(\FREE\) state, where Lemma~\ref{lem:dense-fair-while-on} gives a fresh fair bit, or in the \(\CRIT\) state
while termination has not yet occurred, where the fixed tap forces the output to be \(1\). Thus these nonterminal live
outputs stochastically dominate fair coins in the lower-tail direction, exactly as in
Lemma~\ref{lem:fixedtap-surviving-episode-tail}. Across the \(r\) disjoint intervals there are at least
\((T-rm)_+\) such chargeable live outputs. If the total block weight is at most \(d\), then these chargeable outputs
also contain at most \(d\) ones. The binomial lower-tail bound follows by the same stopping-time induction as
Lemma~\ref{lem:dense-on-time-domination}.
\end{proof}

\begin{corollary}[Full weight-slice paired-cover bound]
\label{cor:fixedtap-weight-slice-paired-cover}
Let \(h\ge 1\), let \(U\) be uniform on the weight-\(h\) slice, let \(d=\lfloor\delta N\rfloor\), and fix
\(1\le L\le N\). Put
\[
\Gamma_{N,h,L}(r,T)
:=
\frac{
\binom{N-T+r}{r}
\binom{T-r-1}{r-1}
\binom{T+L-2r}{h-2r}
}{
\binom{N}{h}
},
\]
with the convention that impossible binomial coefficients are zero. Then the fixed-tap dense recursive inner satisfies
\[
p_h(\delta)
\le
\frac{\binom{L}{h}}{\binom{N}{h}}
\;+\;
\Pr[\Binomial(L,\tfrac12)\le d]
\;+\;
\sum_{r=1}^{\lfloor h/2\rfloor}
\sum_{T=2r}^{\min\{N,r(L+1)\}}
\Gamma_{N,h,L}(r,T)\,
2^{-r(m-1)}
\Pr\!\left[\Binomial((T-rm)_+,\tfrac12)\le d\right].
\]
\end{corollary}

\begin{proof}
If \(E_L(U)=\emptyset\), then all \(h\) input ones lie in the last \(L\) positions, which gives the first term.
Otherwise consider the first ON episode started by an input-one position in \(E_L(U)\). If this episode, or any later
episode started inside \(E_L(U)\), survives for \(L\) output times, then choosing the first such episode predictably and
applying Lemma~\ref{lem:fixedtap-surviving-episode-tail} gives the displayed binomial-tail term.

It remains to bound the event that every episode started inside \(E_L(U)\) terminates within \(L\) output times.
Lemma~\ref{lem:fixedtap-no-survival-paired-cover} then produces a paired \(L\)-cover. For a fixed ordered cover with
\(r\) intervals and total live length \(T\), Lemma~\ref{lem:fixedtap-paired-cover-charge} bounds the joint probability
of realizing that cover and having \(\wt(Y)\le d\) by the stochastic factor displayed above.

We now union bound over possible covers and average over the weight-\(h\) input slice. The number of ordered disjoint
intervals in \([N]\) with lengths at least \(2\) and total length \(T\) is at most
\[
\binom{T-r-1}{r-1}\binom{N-T+r}{r}.
\]
Indeed, first choose the \(r\) lengths \(\ell_a\ge2\) summing to \(T\), then distribute the \(N-T\) unused positions
among the \(r+1\) gaps before, between, and after the intervals. Once the \(2r\) endpoints are forced into the support
of \(U\), every remaining input-one position must lie either in the interiors of the cover intervals or in the last
\(L\) positions; this gives at most \(\binom{T+L-2r}{h-2r}\) possible completions. Dividing by \(\binom{N}{h}\) gives
\(\Gamma_{N,h,L}(r,T)\). The interval count ignores the extra restriction that starts lie in \(E_L(U)\), so it is an
upper bound. Summing over \(r\) and \(T\) proves the claim.
\end{proof}


\subsection{Linear OFF-budget refinement}
\begin{lemma}[Long zero-gap mass is exponentially unlikely]
\label{lem:dense-long-gap-mass}
Fix constants \(0<\eta_0<\eta_1<1\) and \(b>0\). For a binary word \(u\), let \(Z_G(u)\) be the number of zero
coordinates belonging to a zero-run of length greater than \(G\). Then there are constants \(G_0\) and
\(c_{\mathrm{gap}}>0\), depending only on \(\eta_0,\eta_1,b\), such that for every fixed \(G\ge G_0\), uniformly over
\(\eta_0N\le w\le \eta_1N\),
\[
\PP[Z_G(U)\ge bN\mid \wt(U)=w]
\le
2^{-c_{\mathrm{gap}}N}
\]
for all sufficiently large \(N\).
\end{lemma}

\begin{proof}
If \(Z_G(u)\ge bN\), then the zero-runs of length greater than \(G\) contain at least
\[
q:=\left\lfloor \frac{bN}{2G}\right\rfloor
\]
pairwise disjoint all-zero intervals of length \(G\). Indeed, a zero-run of length \(\ell>G\) contains
\(\lfloor \ell/G\rfloor\ge \ell/(2G)\) such intervals.

For a fixed collection of \(q\) disjoint length-\(G\) intervals, the probability that all coordinates in their union
are zero under the uniform weight-\(w\) model is
\[
\frac{\binom{N-qG}{w}}{\binom{N}{w}}
\le
\left(1-\frac{qG}{N}\right)^w
\le
2^{-\eta_0 qG\log_2 e}.
\]
There are at most \(\binom{N}{q}\) possible choices for the left endpoints of such an interval collection, so the union
bound gives
\[
\PP[Z_G(U)\ge bN\mid \wt(U)=w]
\le
\binom{N}{q}2^{-\eta_0 qG\log_2 e}.
\]
For all large \(N\), \(q/N\le b/(2G)\) and \(qG/N\ge b/4\). Hence
\[
\frac1N\log_2 \PP[Z_G(U)\ge bN\mid \wt(U)=w]
\le
H_2\!\left(\frac{b}{2G}\right)-\frac{\eta_0 b}{4}\log_2 e+o(1).
\]
Choosing \(G\) large enough makes the right-hand side negative, proving the claim.
\end{proof}

\begin{lemma}[Short-gap witness packing]
\label{lem:dense-short-gap-witness-packing}
Fix constants \(a>0\) and \(G\ge 1\), and let
\[
m=\gamma\log_2 N
\qquad(\gamma>1).
\]
There is a constant \(c_{\mathrm{pack}}=c_{\mathrm{pack}}(a,G,\gamma)>0\) such that, for every input weight \(w\),
\[
\PP[\mathrm{OffAfter}(T)\ge aN\ \wedge\ Z_G(U)<aN/4\mid \wt(U)=w]
\le
2^{-c_{\mathrm{pack}}N}
\]
for all sufficiently large \(N\).
\end{lemma}

\begin{proof}
Call a post-start OFF episode \emph{long-gap-touching} if it contains a zero coordinate that belongs to a zero-run of
length greater than \(G\). On the event \(Z_G(U)<aN/4\), the number of such long zero coordinates is less than \(aN/4\).
Each long-gap-touching OFF episode contains at least one such coordinate and contributes at most one additional
reactivating \(1\)-coordinate after its zero part. Hence all long-gap-touching OFF episodes together contribute less
than \(aN/2\) OFF times.

Therefore, if also \(\mathrm{OffAfter}(T)\ge aN\), then at least \(aN/2\) OFF times lie in OFF episodes that do not
touch a long zero-run. Such an episode is contained in a zero-run of length at most \(G\), together with possibly one
reactivating \(1\)-coordinate, so it has length at most \(G+1\). Thus the trajectory must contain at least
\[
K:=\left\lfloor \frac{aN}{2(G+1)}\right\rfloor
\]
distinct post-start ON\(\to\)OFF transitions.

Every ON\(\to\)OFF transition has a length-\(m\) all-zero output witness by
Lemma~\ref{lem:dense-off-requires-mzeros}. Witnesses belonging to distinct ON\(\to\)OFF transitions are disjoint: after
one transition, the encoder must first see a \(1\)-output to become ON again, and the next all-zero witness cannot
include that \(1\)-output. Hence \(K\) such transitions require at least \(Km\) distinct time coordinates. But
\[
Km
\ge
\left(\frac{aN}{3(G+1)}\right)\gamma\log_2N
>
N
\]
for all sufficiently large \(N\). Thus the event in the lemma is eventually empty, which is stronger than the claimed
exponential bound.
\end{proof}

\begin{lemma}[Linear-weight OFF-time is exponentially unlikely]
\label{lem:dense-linear-offtime-target}
Fix constants \(\eta_0,\eta_1,a\) with \(0<\eta_0<\eta_1<1\) and \(a>0\), and let
\[
m=\gamma\log_2 N
\qquad(\gamma>1).
\]
There is a constant \(c_{\mathrm{off}}=c_{\mathrm{off}}(\eta_0,\eta_1,a,\gamma)>0\) such that, uniformly for
\[
\eta_0N\le w\le \eta_1N,
\]
\begin{equation}
\label{eq:dense-linear-offtime-target}
\PP[\mathrm{OffAfter}(T)\ge aN\mid \wt(U)=w]
\;\le\;
2^{-c_{\mathrm{off}}N}
\end{equation}
for all sufficiently large \(N\).
\end{lemma}

\begin{proof}
Choose \(G\) large enough for Lemma~\ref{lem:dense-long-gap-mass} with \(b=a/4\). Then there is
\(c_{\mathrm{gap}}>0\) such that
\[
\PP[Z_G(U)\ge aN/4\mid \wt(U)=w]\le 2^{-c_{\mathrm{gap}}N}.
\]
On the complementary event \(Z_G(U)<aN/4\), Lemma~\ref{lem:dense-short-gap-witness-packing} gives
\[
\PP[\mathrm{OffAfter}(T)\ge aN\ \wedge\ Z_G(U)<aN/4\mid \wt(U)=w]\le 2^{-c_{\mathrm{pack}}N}.
\]
Taking \(c_{\mathrm{off}}<\min\{c_{\mathrm{gap}},c_{\mathrm{pack}}\}\) and union bounding proves
\eqref{eq:dense-linear-offtime-target}.
\end{proof}

\begin{corollary}[First-start plus OFF-budget linear tail]
\label{cor:dense-first-start-offbudget-target}
Assume Lemma~\ref{lem:dense-linear-offtime-target}. Fix \(\delta\in(0,1/2)\), an OFF budget \(a>0\), and
\(w=\lfloor\eta N\rfloor\) in the linear window of that lemma. Then
\[
p_w(\delta)
\;\le\;
\sum_{t=1}^{N-w+1}
\frac{\binom{N-t}{w-1}}{\binom{N}{w}}\,
\PP[\Binomial(N-t-\lceil aN\rceil,\tfrac12)\le \lfloor\delta N\rfloor]
\;+\;
2^{-c_{\mathrm{off}}N},
\]
where binomial terms with \(N-t-\lceil aN\rceil<0\) are interpreted as \(1\).
\end{corollary}

\begin{proof}
On the event \(\mathrm{OffAfter}(T)<aN\), the suffix after \(T=t\) contains at least \(N-t-\lceil aN\rceil\) ON times. Take
the first \(N-t-\lceil aN\rceil\) post-start ON times as the predictable stopping times in
Lemma~\ref{lem:dense-on-time-domination}. If \(\wt(Y)\le \lfloor\delta N\rfloor\), then those ON outputs contain at
most \(\lfloor\delta N\rfloor\) ones, so the joint probability is bounded by the displayed binomial lower tail. The
complement is controlled by Lemma~\ref{lem:dense-linear-offtime-target}.
\end{proof}
